\subsection{Variance reduction of the joint estimator}
\label{sec:variance-reduction}

We formalise why modelling the joint distribution
$p(Y_{\mathrm{do}(0)}, Y_{\mathrm{do}(1)} \mid X)$
strictly improves the variance of the CATE estimate compared to modelling
the two marginals independently, whenever the potential outcomes are
positively correlated.

\paragraph{Setup.}
Fix a query covariate $X$ and let the two potential outcomes at $X$ follow
the bivariate normal
\begin{equation}
\begin{pmatrix} Y_{\mathrm{do}(0)} \\ Y_{\mathrm{do}(1)} \end{pmatrix}
\;\sim\;
\mathcal{N}\!\left(\!
\begin{pmatrix} \mu_0 \\ \mu_1 \end{pmatrix},\;
\begin{pmatrix}
\sigma_0^2 & \rho\,\sigma_0\sigma_1 \\
\rho\,\sigma_0\sigma_1 & \sigma_1^2
\end{pmatrix}\!
\right),
\qquad \rho \in [-1, 1].
\label{eq:bvn-setup}
\end{equation}
The true individual treatment effect is
$\tau \;=\; Y_{\mathrm{do}(1)} - Y_{\mathrm{do}(0)}$,
so its \emph{unit-level} variance is
\begin{equation}
\mathrm{Var}(\tau \mid X) \;=\; \sigma_0^2 + \sigma_1^2 - 2\rho\,\sigma_0\sigma_1.
\label{eq:var-tau}
\end{equation}
Assume $N$ context units are sampled from $p(X)\!\cdot\!p(T,Y\mid X)$ with
$N/2$ assigned to each treatment arm ($T\in\{0,1\}$).

\paragraph{Estimator 1 --- Two independent marginals (UWYK-style).}
A marginal-only model estimates $\mu_0, \mu_1$ by two independent arm-specific
sample means and takes the difference:
\begin{equation}
\hat{\tau}^{\,\mathrm{marg}}
\;=\;
\underbrace{\tfrac{1}{N/2}\!\!\sum_{i:\,T_i=1}\!Y_i}_{\hat{\mu}_1}
\;-\;
\underbrace{\tfrac{1}{N/2}\!\!\sum_{i:\,T_i=0}\!Y_i}_{\hat{\mu}_0}.
\end{equation}
The two summands are independent random variables, so
\begin{equation}
\mathrm{Var}\!\left(\hat{\tau}^{\,\mathrm{marg}}\right)
\;=\;
\frac{\sigma_0^2}{N/2} + \frac{\sigma_1^2}{N/2}
\;=\;
\frac{2\,(\sigma_0^2 + \sigma_1^2)}{N}.
\label{eq:var-marg}
\end{equation}
The correlation $\rho$ does not appear: a marginal-only estimator cannot
exploit within-unit dependence.

\paragraph{Estimator 2 --- Joint (Ours).}
A joint model represents the paired distribution
$p(Y_{\mathrm{do}(0)}, Y_{\mathrm{do}(1)} \mid X)$ and can, in the limit of an
expressive model, produce paired predictions
$(\hat{Y}_{\mathrm{do}(0)}(i), \hat{Y}_{\mathrm{do}(1)}(i))$
per context unit. The corresponding estimator is
\begin{equation}
\hat{\tau}^{\,\mathrm{joint}}
\;=\;
\frac{1}{N}\sum_{i=1}^{N}
\Bigl(\hat{Y}_{\mathrm{do}(1)}(i) - \hat{Y}_{\mathrm{do}(0)}(i)\Bigr).
\end{equation}
Its variance inherits the joint's coupling term:
\begin{equation}
\mathrm{Var}\!\left(\hat{\tau}^{\,\mathrm{joint}}\right)
\;=\;
\frac{\sigma_0^2 + \sigma_1^2 - 2\rho\,\sigma_0\sigma_1}{N}.
\label{eq:var-joint}
\end{equation}

\begin{theorem}[Variance-ratio guarantee]
\label{thm:variance-ratio}
Under the setup of \eqref{eq:bvn-setup}, the marginal-to-joint variance
ratio equals
\begin{equation}
R(\rho) \;\coloneqq\;
\frac{\mathrm{Var}\!\left(\hat{\tau}^{\,\mathrm{marg}}\right)}
{\mathrm{Var}\!\left(\hat{\tau}^{\,\mathrm{joint}}\right)}
\;=\;
\frac{2\,(\sigma_0^2 + \sigma_1^2)}{\sigma_0^2 + \sigma_1^2 - 2\rho\,\sigma_0\sigma_1}
\;\;\;\underset{\sigma_0=\sigma_1}{=}\;\;\;
\frac{2}{1-\rho}.
\label{eq:ratio}
\end{equation}
In particular $R$ is strictly increasing in $\rho$ on $[0,1)$ with
$R(0)=2$ and $\lim_{\rho\to 1^-}R(\rho)=\infty$.
The joint estimator strictly dominates the marginal estimator in
mean-square error whenever $\rho > 0$, and the improvement is unbounded as
$\rho \to 1$.
\end{theorem}

\begin{proof}
Substitute \eqref{eq:var-marg} and \eqref{eq:var-joint} into the definition
of $R(\rho)$. Setting $\sigma_0 = \sigma_1 = \sigma$ simplifies the ratio
to $\frac{4\sigma^2}{2\sigma^2(1-\rho)}=\frac{2}{1-\rho}$.
Both estimators are unbiased for $\tau$, so lower variance implies lower MSE.
\qed
\end{proof}

\paragraph{Empirical corollary (Figure~\ref{fig:variance-reduction}).}
The Monte-Carlo simulation in Figure~\ref{fig:variance-reduction} uses
$\sigma_0=\sigma_1=1$ and $N=50$. Theorem~\ref{thm:variance-ratio} predicts
\[
R(0) = 2, \qquad R(0.5) = 4, \qquad R(0.9) = 20,
\]
matching the observed empirical ratios of $2.0\times$, $4.1\times$, and
$19.9\times$ across $10{,}000$ trials to within Monte-Carlo error. This
holds under any convex loss (squared error, in particular), so the joint
model dominates the marginal-only baseline irrespective of whether the
point estimate is the posterior mean or mode.
